\section{THIẾT KẾ HỆ THỐNG MẠNG}

\subsection{Mô hình thiết kế tổng thể}
Dựa trên yêu cầu của đề tài, mạng Campus 3 lớp được thiết kế trên môi trường Mininet với sự phân tách rõ ràng về mặt chức năng giữa các lớp Core, Distribution, Access và vùng an ninh DMZ. 

% CHỤP ẢNH MINH CHỨNG:
% Yêu cầu: Chụp sơ đồ Topology đã được tạo bằng file draw_topology.py (nằm trong thư mục docs/network_topology.png).
% Lưu ảnh vào thư mục image/ với tên topology.png và bỏ comment đoạn code bên dưới để hiển thị ảnh.
% \begin{figure}[H]
%     \centering
%     \includegraphics[width=1\textwidth]{image/topology.png}
%     \caption{Sơ đồ Topology mạng Campus 3 lớp và DMZ}
%     \label{fig:topology}
% \end{figure}

Sơ đồ mô phỏng bao gồm:
\begin{itemize}
    \item \textbf{Khối Biên (Edge) và Tường lửa (Firewall):} Router \texttt{fw} đóng vai trò kết nối với mạng External (gia lập Internet) và kiểm soát an ninh mạng thông qua iptables.
    \item \textbf{Vùng DMZ:} Gồm 1 switch (\texttt{s1}) và 2 máy chủ (\texttt{web1}, \texttt{web2}) chạy dịch vụ HTTP và iPerf.
    \item \textbf{Lớp Core:} Router \texttt{core} là bộ định tuyến trung tâm kết nối Tường lửa với các khối Distribution.
    \item \textbf{Lớp Distribution:} Bao gồm các Router \texttt{dist1} và \texttt{dist2}, phục vụ như Default Gateway cho các mạng LAN bên dưới.
    \item \textbf{Lớp Access:} Bao gồm các Switch \texttt{s2}, \texttt{s3} kết nối trực tiếp các Host nội bộ (\texttt{h1}, \texttt{h2}, \texttt{h3}, \texttt{h4}).
\end{itemize}

\subsection{Thiết kế quy hoạch địa chỉ IP và Định tuyến}

\subsubsection{Bảng quy hoạch IP}
Địa chỉ IP được chia thành các phân đoạn mạng (Subnet) phù hợp với từng khu vực:
\begin{table}[H]
    \centering
    \caption{Quy hoạch địa chỉ IP cho các phân đoạn mạng}
    \begin{tabular}{|l|l|l|}
    \hline
    \rowcolor[HTML]{EFEFEF} 
    \textbf{Khu vực} & \textbf{Subnet} & \textbf{Ghi chú / Thiết bị} \\ \hline
    Internet (External WAN) & 203.0.113.0/30 & ext=.1, fw=.2 \\ \hline
    Kết nối Firewall - Core & 172.16.0.0/30 & fw=.1, core=.2 \\ \hline
    Vùng DMZ & 10.10.10.0/24 & fw=.1, web1=.11, web2=.12 \\ \hline
    Core - Dist1 & 172.16.1.0/30 & core=.1, dist1=.2 \\ \hline
    Core - Dist2 & 172.16.2.0/30 & core=.1, dist2=.2 \\ \hline
    LAN Access 1 & 192.168.10.0/24 & dist1=.1, h1=.11, h2=.12 \\ \hline
    LAN Access 2 & 192.168.20.0/24 & dist2=.1, h3=.11, h4=.12 \\ \hline
    \end{tabular}
\end{table}

\subsubsection{Cấu hình định tuyến OSPF}
Giao thức OSPF được cấu hình trên các Router \texttt{fw}, \texttt{core}, \texttt{dist1}, \texttt{dist2} chạy chung trên Area 0. Các nút sử dụng phần mềm định tuyến mã nguồn mở FRRouting (FRR) để mô phỏng Router thật.
Mỗi Router được gán một Router-ID thông qua giao diện Loopback ảo: \texttt{10.255.1.10} cho \texttt{fw}, \texttt{10.255.1.1} cho \texttt{core}, \texttt{10.255.1.2} và \texttt{10.255.1.3} cho \texttt{dist1, dist2}.

\subsection{Thiết kế chính sách NAT và PAT}
Tường lửa \texttt{fw} thực hiện chức năng biên dịch địa chỉ nhằm che giấu dải IP Private và cung cấp dịch vụ công cộng cho máy chủ nội bộ.

\begin{itemize}
    \item \textbf{PAT Overload (MASQUERADE):} Toàn bộ lưu lượng từ mạng \texttt{192.168.10.0/24} và \texttt{192.168.20.0/24} đi qua cổng \texttt{fw-eth0} ra Internet sẽ được biên dịch (MASQUERADE) sang địa chỉ Public của Firewall \texttt{203.0.113.2}.
    \item \textbf{Static NAT:} Gán địa chỉ IP Public ảo (IP Alias) trên cổng WAN của Firewall để dịch tĩnh vào địa chỉ DMZ Private:
    \begin{itemize}
        \item \texttt{100.0.0.11} $\leftrightarrow$ \texttt{10.10.10.11} (\texttt{web1})
        \item \texttt{100.0.0.12} $\leftrightarrow$ \texttt{10.10.10.12} (\texttt{web2})
    \end{itemize}
\end{itemize}

\subsection{Thiết kế kiểm soát truy cập (ACL)}
Bảo mật hệ thống được thiết kế theo chính sách "Default DROP" ở chuỗi FORWARD của Firewall \texttt{fw}. Chỉ những luồng dữ liệu hợp lệ mới được cấp phép:

\begin{enumerate}
    \item \textbf{Cho phép kết nối đã thiết lập:} Duy trì các phiên làm việc (state ESTABLISHED, RELATED).
    \item \textbf{Extended ACL:} Cho phép các dải mạng nội bộ (\texttt{192.168.10.0/24} và \texttt{192.168.20.0/24}) truy cập cổng HTTP(80) và HTTPS(443) trên máy chủ DMZ. Các truy cập khác bị cấm.
    \item \textbf{Standard ACL chặn SSH:} Tường lửa chặn riêng rẽ các kết nối giao thức TCP cổng 22 xuất phát từ mạng \texttt{192.168.20.0/24} vào DMZ.
    \item \textbf{Ngăn chặn rủi ro từ DMZ:} Luồng dữ liệu khởi tạo từ DMZ \texttt{10.10.10.0/24} đi vào các mạng nội bộ \texttt{192.168.0.0/16} hoàn toàn bị DROP để tránh DMZ trở thành bàn đạp tấn công (Pivot Attack).
    \item \textbf{Cho phép External Ping và Web:} Người dùng từ mạng bên ngoài (\texttt{ext}) được phép truy cập cổng 80, 443 và ping (ICMP) đến IP Static NAT \texttt{10.10.10.0/24} trên DMZ.
\end{enumerate}

\subsection{Thuật toán Cân bằng tải tự động}
Nhằm hạn chế tình trạng thắt nút cổ chai (Bottleneck) khi lưu lượng vào \texttt{web1} quá lớn, một script Python giám sát nền (\texttt{load\_balancer.py}) được tích hợp. Script hoạt động dựa trên các ngưỡng cấu hình:
\begin{itemize}
    \item \textbf{Trạng thái Normal (Bình thường):} \texttt{web1} chịu 100\% lưu lượng.
    \item \textbf{Threshold 80\%:} Nếu script phát hiện băng thông tại \texttt{web1} vượt quá 80\% giới hạn, lập tức thực thi chuỗi lệnh iptables trên Tường lửa để \texttt{REDIRECT/DNAT} kết nối mới từ Internet sang IP \texttt{10.10.10.12} (\texttt{web2}).
    \item \textbf{Threshold 20\%:} Khi tải tổng giảm về mức thấp, Tường lửa xóa quy tắc REDIRECT, đưa lưu lượng định tuyến trở lại \texttt{web1}.
\end{itemize}
